% GENERATED FILE. Edit canonical lessons, then run npm run generate:books.
% canonical-source-hash: fnv1a64:2a9eb80105c62069
% canonical-lessons: TA-C416-patakotti, TA-C416-tutuppu, TA-C416-toni, TA-C416-kattumaram, TA-C416-valai

\chapter{Boats and Oars}
\label{ch:ta-boats-and-oars}
\hlchaptermodality{416}

\begin{chapteropening}
\textbf{By the end of this chapter:} \emph{I can talk about a boatman, oars, small boats, a catamaran and a fishing net.}

Complete the last lesson of chapter 416: I can talk about a boatman, oars, small boats, a catamaran and a fishing net.
\end{chapteropening}

\section[\texorpdfstring{\ta{படகோட்டி} (paṭakōṭṭi)}{படகோட்டி (paṭakōṭṭi)}]{\ta{படகோட்டி} (paṭakōṭṭi) — a boatman}

\label{lesson:TA-C416-patakotti}



\begin{quote}
Before the new one: say the Tamil for an acknowledgement, then the Tamil for to arrange.
\end{quote}

\subsection*{You'll want to know: \ta{படகோட்டி}}
\textbf{\ta{படகோட்டி}} — \emph{paṭakōṭṭi} — \textquotedblleft{}a boatman\textquotedblright{}.

\begin{quote}
\ta{படகோட்டி} \ta{எங்களைத்} \ta{தீவுக்குக்} \ta{கொண்டு} \ta{போனார்}.
\end{quote}

\begin{quote}
\emph{paṭakōṭṭi eṅkaḷait tīvukkuk koṇṭu pōṉār.}
\end{quote}

\begin{quote}
The boatman took us to the island.
\end{quote}

\ta{படகு}, a boat, + \ta{ஓட்டி}, one who drives.

\begin{cousinweb}
\textbf{\ta{படகோட்டி}} (\emph{paṭakōṭṭi}) is built from \ta{படகு}, \textquotedblleft{}a boat\textquotedblright{} (DEDR 3838), and \ta{ஓட்டி}, \textquotedblleft{}one who drives\textquotedblright{}, from \ta{ஓட்டு} (DEDR 1041).
\end{cousinweb}

\begin{grammarlens}[title={What you've built}]
One more word for getting around.
\end{grammarlens}

\subsection*{Your turn}
\begin{itemize}
\raggedright
  \item \emph{Say it:} \emph{paṭakōṭṭi}
  \item \emph{Say it:} \emph{paṭakōṭṭi}, once more
  \item \emph{Say it:} \emph{paṭakōṭṭi eṅkaḷait tīvukkuk koṇṭu pōṉār.} — The boatman took us to the island.
  \item \emph{Recall:} say the Tamil for an acknowledgement, then the Tamil for to arrange, then say \emph{paṭakōṭṭi} again
\end{itemize}

\begin{tcolorbox}[breakable,colback=teal!4,colframe=teal!35!black,title={Before you move on}]
What does \ta{படகோட்டி} mean? (\textquotedblleft{}A boatman\textquotedblright{}.) Say it once more.
\end{tcolorbox}
\section[\texorpdfstring{\ta{துடுப்பு} (tuṭuppu)}{துடுப்பு (tuṭuppu)}]{\ta{துடுப்பு} (tuṭuppu) — an oar; a paddle}

\label{lesson:TA-C416-tutuppu}



\begin{quote}
Before the new one: say the Tamil for to arrange, then the Tamil for a boatman.
\end{quote}

\subsection*{You'll want to know: \ta{துடுப்பு}}
\textbf{\ta{துடுப்பு}} — \emph{tuṭuppu} — \textquotedblleft{}an oar; a paddle\textquotedblright{}.

\begin{quote}
\ta{துடுப்பு} \ta{இல்லாமல்} \ta{படகு} \ta{நகராது}.
\end{quote}

\begin{quote}
\emph{tuṭuppu illāmal paṭaku nakarātu.}
\end{quote}

\begin{quote}
Without an oar the boat will not move.
\end{quote}

Also a ladle or a flat spoon for stirring.

\begin{cousinweb}
\textbf{\ta{துடுப்பு}} (\emph{tuṭuppu}) is native Dravidian (DEDR 3299), \textquotedblleft{}a ladle, a spatula, an oar\textquotedblright{}.
\end{cousinweb}

\begin{grammarlens}[title={What you've built}]
One more word for getting around.
\end{grammarlens}

\subsection*{Your turn}
\begin{itemize}
\raggedright
  \item \emph{Say it:} \emph{tuṭuppu}
  \item \emph{Say it:} \emph{tuṭuppu}, once more
  \item \emph{Say it:} \emph{tuṭuppu illāmal paṭaku nakarātu.} — Without an oar the boat will not move.
  \item \emph{Recall:} say the Tamil for to arrange, then the Tamil for a boatman, then say \emph{tuṭuppu} again
\end{itemize}

\begin{tcolorbox}[breakable,colback=teal!4,colframe=teal!35!black,title={Before you move on}]
What does \ta{துடுப்பு} mean? (\textquotedblleft{}An oar; a paddle\textquotedblright{}.) Say it once more.
\end{tcolorbox}
\section[\texorpdfstring{\ta{தோணி} (tōṇi)}{தோணி (tōṇi)}]{\ta{தோணி} (tōṇi) — a dinghy; a small boat}

\label{lesson:TA-C416-toni}



\begin{quote}
Before the new one: say the Tamil for a boatman, then the Tamil for an oar.
\end{quote}

\subsection*{You'll want to know: \ta{தோணி}}
\textbf{\ta{தோணி}} — \emph{tōṇi} — \textquotedblleft{}a dinghy; a small boat\textquotedblright{}.

\begin{quote}
\ta{மீனவர்கள்} \ta{தோணியில்} \ta{கடலுக்குப்} \ta{போனார்கள்}.
\end{quote}

\begin{quote}
\emph{mīṉavarkaḷ tōṇiyil kaṭalukkup pōṉārkaḷ.}
\end{quote}

\begin{quote}
The fishermen went to sea in a small boat.
\end{quote}

The plural is \ta{தோணிகள்}.

\begin{cousinweb}
\textbf{\ta{தோணி}} (\emph{tōṇi}) is usually traced to Sanskrit \emph{droṇī}, \textquotedblleft{}a wooden trough, a tub\textquotedblright{} (CDIAL 6641).
\end{cousinweb}

\begin{grammarlens}[title={What you've built}]
One more word for getting around.
\end{grammarlens}

\subsection*{Your turn}
\begin{itemize}
\raggedright
  \item \emph{Say it:} \emph{tōṇi}
  \item \emph{Say it:} \emph{tōṇi}, once more
  \item \emph{Say it:} \emph{mīṉavarkaḷ tōṇiyil kaṭalukkup pōṉārkaḷ.} — The fishermen went to sea in a small boat.
  \item \emph{Recall:} say the Tamil for a boatman, then the Tamil for an oar, then say \emph{tōṇi} again
\end{itemize}

\begin{tcolorbox}[breakable,colback=teal!4,colframe=teal!35!black,title={Before you move on}]
What does \ta{தோணி} mean? (\textquotedblleft{}A dinghy; a small boat\textquotedblright{}.) Say it once more.
\end{tcolorbox}
\section[\texorpdfstring{\ta{கட்டுமரம்} (kaṭṭumaram)}{கட்டுமரம் (kaṭṭumaram)}]{\ta{கட்டுமரம்} (kaṭṭumaram) — a catamaran; a log raft}

\label{lesson:TA-C416-kattumaram}



\begin{quote}
Before the new one: say the Tamil for an oar, then the Tamil for a dinghy.
\end{quote}

\subsection*{You'll want to know: \ta{கட்டுமரம்}}
\textbf{\ta{கட்டுமரம்}} — \emph{kaṭṭumaram} — \textquotedblleft{}a catamaran; a log raft\textquotedblright{}.

\begin{quote}
\ta{கட்டுமரம்} \ta{அலைகளின்} \ta{மேல்} \ta{மிதந்தது}.
\end{quote}

\begin{quote}
\emph{kaṭṭumaram alaikaḷiṉ mēl mitantatu.}
\end{quote}

\begin{quote}
The log raft floated on the waves.
\end{quote}

English catamaran comes from this word.

\begin{cousinweb}
\textbf{\ta{கட்டுமரம்}} (\emph{kaṭṭumaram}) is built from \ta{கட்டு}, \textquotedblleft{}to tie\textquotedblright{} (DEDR 1147), and \ta{மரம்}, \textquotedblleft{}a tree, wood\textquotedblright{} (DEDR 4711): logs tied together.
\end{cousinweb}

\begin{grammarlens}[title={What you've built}]
One more word for getting around.
\end{grammarlens}

\subsection*{Your turn}
\begin{itemize}
\raggedright
  \item \emph{Say it:} \emph{kaṭṭumaram}
  \item \emph{Say it:} \emph{kaṭṭumaram}, once more
  \item \emph{Say it:} \emph{kaṭṭumaram alaikaḷiṉ mēl mitantatu.} — The log raft floated on the waves.
  \item \emph{Recall:} say the Tamil for an oar, then the Tamil for a dinghy, then say \emph{kaṭṭumaram} again
\end{itemize}

\begin{tcolorbox}[breakable,colback=teal!4,colframe=teal!35!black,title={Before you move on}]
What does \ta{கட்டுமரம்} mean? (\textquotedblleft{}A catamaran; a log raft\textquotedblright{}.) Say it once more.
\end{tcolorbox}
\section[\texorpdfstring{\ta{வலை} (valai)}{வலை (valai)}]{\ta{வலை} (valai) — a net}

\label{lesson:TA-C416-valai}



\begin{quote}
Before the new one: say the Tamil for a dinghy, then the Tamil for a catamaran.
\end{quote}

\subsection*{You'll want to know: \ta{வலை}}
\textbf{\ta{வலை}} — \emph{valai} — \textquotedblleft{}a net\textquotedblright{}.

\begin{quote}
\ta{மீனவர்கள்} \ta{கடலில்} \ta{வலை} \ta{வீசினார்கள்}.
\end{quote}

\begin{quote}
\emph{mīṉavarkaḷ kaṭalil valai vīciṉārkaḷ.}
\end{quote}

\begin{quote}
The fishermen cast a net into the sea.
\end{quote}

\ta{வலை} \ta{வீசு} is to cast a net; \ta{கொசுவலை}, a mosquito net.

\begin{cousinweb}
\textbf{\ta{வலை}} (\emph{valai}) is native Dravidian (DEDR 5288), \textquotedblleft{}a net\textquotedblright{}.
\end{cousinweb}

\begin{grammarlens}[title={What you've built}]
One more word for getting around.
\end{grammarlens}

\subsection*{Your turn}
\begin{itemize}
\raggedright
  \item \emph{Say it:} \emph{valai}
  \item \emph{Say it:} \emph{valai}, once more
  \item \emph{Say it:} \emph{mīṉavarkaḷ kaṭalil valai vīciṉārkaḷ.} — The fishermen cast a net into the sea.
  \item \emph{Recall:} say the Tamil for a dinghy, then the Tamil for a catamaran, then say \emph{valai} again
\end{itemize}

\begin{tcolorbox}[breakable,colback=teal!4,colframe=teal!35!black,title={Before you move on}]
What does \ta{வலை} mean? (\textquotedblleft{}A net\textquotedblright{}.) Say it once more.
\end{tcolorbox}
